\chapter[Chapter \thechapter]{}
\lettrine{T}{he} next day dawned bright and fair, and Wimsey felt a certain exhilaration as he purred down to Tweedling Parva. <Mrs~Merdle> the car, so called because, like that celebrated lady, she was averse to <row,> was sparking merrily on all twelve cylinders, and there was a touch of frost in the air. These things conduce to high spirits.

Wimsey reached his destination about 10 o'clock, and was directed to the vicarage, one of those large, rambling and unnecessary structures which swallow the incumbent's income during his life and land his survivors with a heavy bill for dilapidations as soon as he is dead.

The Rev. Arthur Boyes was at home, and would be happy to see Lord~Peter Wimsey.

The clergyman was a tall, faded man, with lines of worry deeply engraved upon his face, and mild blue eyes a little bewildered by the disappointing difficulty of things in general. His black coat was old, and hung in depressed folds from his stooping, narrow shoulders. He gave Wimsey a thin hand and begged him to be seated.

Lord~Peter found it a little difficult to explain his errand. His name evidently aroused no associations in the mind of this gentle and unworldly parson. He decided not to mention his hobby of criminal investigation, but to represent himself, with equal truth, as a friend of the prisoner's. That might be painful, but it would be at least intelligible. Accordingly, he began, with some hesitation:

<I'm fearfully sorry to trouble you, especially as it's all so very distressin' and all that, but it's about the death of your son, and the trial and so on. Please don't think I'm wanting to make an interfering nuisance of myself, but I'm deeply interested—personally interested. You see, I know Miss~Vane—I—in fact I like her very much, don't you know, and I can't help thinking there's a mistake somewhere and—and I should like to get it put right if possible.>

<Oh—oh, yes!> said Mr~Boyes. He carefully polished a pair of pince-nez and balanced them on his nose, where they sat crookedly. He peered at Wimsey and seemed not to dislike what he saw, for he went on:

<Poor misguided girl! I assure you, I have no vindictive feelings—that is to say, nobody would be more happy than myself to know that she was innocent of this dreadful thing. Indeed, Lord~Peter, even if she were guilty, it would give me great pain to see her suffer the penalty. Whatever we do, we cannot bring back the dead to life, and one would infinitely prefer to leave all vengeance in the hand of Him to whom it belongs. Certainly, nothing could be more terrible than to take the life of an innocent person. It would haunt me to the end of my days if I thought there were the least likelihood of it. And I confess that, when I saw Miss~Vane in court, I had grievous doubts whether the police had done rightly in accusing her.>

<Thank you,> said Wimsey, <it is very kind of you to say that. It makes the job much easier. Excuse me, you say, <when you saw her in court.> You hadn't met her previously?>

<No. I knew, of course, that my unhappy son had formed an illicit connection with a young woman, but—I could not bring myself to see her—and indeed, I believe that she, with very proper feeling, refused to allow Philip to bring her into contact with any of his relations. Lord~Peter, you are a younger man than I am, you belong to my son's generation, and you will perhaps understand that—though he was not bad, not depraved, I will never think that—yet somehow there was not that full confidence between us which there should be between father and son. No doubt I was much to blame. If only his mother had lived\longdash>

<My dear sir,> mumbled Wimsey, <I perfectly understand. It often happens. In fact, it's continually happening. The post-war generation and so on. Lots of people go off the rails a bit—no real harm in 'em at all. Just can't see eye to eye with the older people. It generally wears off in time. Nobody really to blame. Wild oats and, er, all that sort of thing.>

<I could not approve,> said Mr~Boyes, sadly, <of ideas so opposed to religion and morality—perhaps I spoke my mind too openly. If I had sympathised more\longdash>

<It can't be done,> said Wimsey. <People have to work it out for themselves. And, when they write books and so on, and get into that set of people, they tend to express themselves rather noisily, if you see what I mean.>

<Maybe, maybe. But I reproach myself. Still, this does not help you at all. Forgive me. If there is any mistake and the jury were evidently not satisfied, we must use all our endeavours to put it right. How can I assist?>

<Well, first of all,> said Wimsey, <and I'm afraid this is rather a hateful question, did your son ever say anything, or write anything to you which might lead you to think that he—was tired of his life or anything of that kind? I'm sorry.>

<No, no—not at all. I was, of course, asked the same question by the police and by the counsel for the defence. I can truly say that such an idea never occurred to me. There was nothing at all to suggest it.>

<Not even when he parted company with Miss~Vane?>

<Not even then. In fact, I gathered that he was rather more angry than despondent. I must say that it was a surprise to me to hear that, after all that had passed between them, she was unwilling to marry him. I still fail to comprehend it. Her refusal must have come as a great shock to him. He wrote so cheerfully to me about it beforehand. Perhaps you remember the letter?> He fumbled in an untidy drawer. <I have it here, if you would like to look at it.>

<If you would just read the passage, sir,> suggested Wimsey.

<Yes, oh, certainly. Let me see. Yes. <Your morality will be pleased to hear, Dad, that I have determined to regularise the situation, as the good people say.> He had a careless way of speaking and writing sometimes, poor boy, which doesn't do justice to his good heart. Dear me. Yes. <My young woman is a good little soul, and I have made up my mind to do the thing properly. She really deserves it, and I hope that when everything is made respectable, you will extend your paternal recognition to her. I won't ask you to officiate—as you know, the registrar's office is more in my line, and though she was brought up in the odour of sanctity, like myself, I don't think she will insist on the Voice that Breathed o'er Eden. I will let you know when it's to be, so that you can come and give us your blessing (quâ father if not quâ parson) if you should feel so disposed.> You see, Lord~Peter, he quite meant to do the right thing, and I was touched that he should wish for my presence.>

<Quite so,> said Lord~Peter, and thought, <If only that young man were alive, how dearly I should love to kick his bottom for him.>

<Well, then there is another letter, saying that the marriage had fallen through. Here it is. <Dear Dad—sorry, but I'm afraid your congratulations must be returned with thanks. The wedding is off, and the bride has run away. There's no need to go into the story. Harriet has succeeded in making a fool of herself and me, so there's no more to be said.> Then later I heard that he had not been feeling well—but all that you know already.>

<Did he suggest any reason for these illnesses of his?>

<Oh, no—we took it for granted that it was a recurrence of the old gastric trouble. He was never a very robust lad. He wrote in very hopeful mood from Harlech, saying that he was much better, and mentioning his plan of a voyage to Barbados.>

<He did?>

<Yes. I thought it would do him a great deal of good, and take his mind off other things. He spoke of it only as a vague project, not as though anything were settled.>

<Did he say anything more about Miss~Vane?>

<He never mentioned her name to me again until he lay dying.>

<Yes—and what did you think of what he said then?>

<I didn't know what to think. We had no idea of any poisoning then, naturally, and I fancied it must refer to the quarrel between them that had caused the separation.>

<I see. Well now, Mr~Boyes. Supposing it was not self-destruction\longdash>

<I really do not think it could have been.>

<Now is there anybody else at all who could have an interest in his death?>

<Who could there be?>

<No—no other woman, for instance?>

<I never heard of any. And I think I should have done. He was not secretive about these things, Lord~Peter. He was remarkably open and straightforward.>

<Yes,> commented Wimsey internally, <liked to swagger about it, I suppose. Anything to give pain. Damn the fellow.> Aloud he merely said: <There are other possibilities. Did he, for instance, make a will?>

<He did. Not that he had much to leave, poor boy. His books were very cleverly written—he had a fine intellect, Lord~Peter—but they did not bring him in any great sums of money. I helped him with a little allowance, and he managed on that and on what he made from his articles in the periodicals.>

<He left his copyrights to somebody, though, I take it?>

<Yes. He wished to leave them to me, but I was obliged to tell him that I could not accept the bequest. You see, I did not approve of his opinions, and I should not have thought it right to profit by them. No; he left them to his friend Mr~Vaughan.>

<Oh!—may I ask when this will was made?>

<It is dated at the period of his visit to Wales. I believe that before that he had made one leaving everything to Miss~Vane.>

<Indeed!> said Wimsey. <I suppose she knew about it.> His mind reviewed a number of contradictory possibilities, and he added: <But it would not amount to an important sum, in any case?>

<Oh, no. If my son made \textsterling 50 a year by his books, that was the utmost. Though they tell me,> added the old gentleman, with a sad smile, <that, after this, his new book will do better.>

<Very likely,> said Wimsey. <Provided you get into the papers, the delightful reading public don't mind what it's for. Still—Well, that's that. I gather he would have no private money to leave?>

<Nothing whatever. There has never been any money in our family, Lord~Peter, nor yet in my wife's. We're quite the proverbial Church mice.> He smiled faintly at this little clerical jest. <Except, I suppose, for Cremorna Garden.>

<For—I beg your pardon?>

<My wife's aunt, the notorious Cremorna Garden of the 'sixties.>

<Good lord, yes—the actress?>

<Yes. But she, of course, was never, never mentioned. One did not enquire into the way she got her money. No worse than others, I dare say—but in those days we were very easily shocked. We have seen and heard nothing of her for well over fifty years. I believe she is quite childish now.>

<By jove! I'd no idea she was still alive!>

<Yes, I believe she is, though she must be well over ninety. Certainly Philip never had any money from her.>

<Well, that rules money out. Was your son's life insured, by any chance?>

<Not that I ever heard of. We found no policy among his papers, and so far as I know, nobody has made any claim.>

<He left no debts?>

<Only trifling ones—tradesmen's accounts and so on. Perhaps fifty pounds' worth altogether.>

<Thank you so much,> said Wimsey, rising, <that has cleared the ground a good deal.>

<I am afraid it has not got you much farther.>

<It tells me where not to look, at any rate,> said Wimsey, <and that all saves time, you know. It's frightfully decent of you to be bothered with me.>

<Not at all. Ask me anything you want to know. Nobody would be more glad than myself to see that unfortunate young woman cleared.>

Wimsey again thanked him and took his leave. He was a mile up the road before a regretful thought overtook him. He turned Mrs~Merdle's bonnet round, skimmed back to the church, stuffed a handful of treasury notes with some difficulty into the mouth of a box labelled <Church expenses,> and resumed his way to town.

\divider

As he manoeuvred the car through the City, a thought struck him, and instead of heading for Piccadilly, where he lived, he turned off into a street south of the Strand, in which was situated the establishment of Messrs. Grimsby \& Cole, who published the works of Mr~Philip Boyes. After a little delay, he was shown into Mr~Cole's office.

Mr~Cole was a stout and cheerful person, and was much interested to hear that the notorious Lord~Peter Wimsey was concerning himself with the affairs of the equally notorious Mr~Boyes. Wimsey represented that, as a collector of First Editions, he would be glad to secure copies of all Philip Boyes' works. Mr~Cole regretted extremely that he could not help him, and, under the influence of an expensive cigar, became quite confidential.

<Without wishing to seem callous, my dear Lord~Peter,> he said, throwing himself back in his chair, and creasing his three chins into six or seven as he did so, <between you and me, Mr~Boyes could not have done better for himself than to go and get murdered like this. Every copy was sold out a week after the result of the exhumation became known, two large editions of his last book were disposed of before the trial came on—at the original price of \oldmoney{0}{7}{6}, and the libraries clamoured so for the early volumes that we had to reprint the lot. Unfortunately we had not kept the type standing, and the printers had to work night and day, but we did it. We are rushing the three-and-sixpennies through the binders' now, and the shilling edition is arranged for. Positively, I don't think you could get a First Edition in London for love or money. We have nothing here but our own file copies, but we are putting out a special memorial edition, with portraits, on hand-made paper, limited and numbered, at a guinea. Not the same thing of course, but\longdash>

Wimsey begged to put his name down for a set at a guinea a-piece, adding:

<Sad and all that, don't you know, that the author can't benefit by it, what?>

<Deeply distressing,> agreed Mr~Cole, compressing his fat cheeks by two longitudinal folds from the nostril to the mouth. <And sadder still that there can be no more work to come from him. A very talented young man, Lord~Peter. We shall always feel a melancholy pride, Mr~Grimsby and myself, in knowing that we recognised his quality, before there was any likelihood of financial remuneration. A \textit{succès d'estime}, that was all, until this very grievous occurrence. But when the work is good, it is not our habit to boggle about monetary returns.>

<Ah, well!> said Wimsey, <it sometimes pays to cast your bread upon the waters. Quite religious, isn't it—you know, the bit about <plenteously bringing out good works may of thee be plenteously rewarded.> Twenty-fifth after Trinity.>

<Quite,> said Mr~Cole, with a certain lack of enthusiasm, possibly because he was imperfectly acquainted with the book of Common Prayer, or possibly because he detected a hint of mockery in the other's tone. <Well, I have very much enjoyed this chat. I am sorry I can do nothing for you about First Editions.>

Wimsey begged him not to mention it, and with a cordial farewell ran hastily down the stairs.

His next visit was to the office of Mr~Challoner, Harriet Vane's agent. Challoner was an abrupt, dark, militant-looking little man, with untidy hair and thick spectacles.

<Boom?> said he, when Wimsey had introduced himself and mentioned his interest in Miss~Vane. <Yes, of course there is a boom. Rather disgusting, really, but one can't help that. We have to do our best for our client, whatever the circumstances. Miss~Vane's books have always sold reasonably well—round about the three or four thousand mark in this country—but of course this business has stimulated things enormously. The last book has gone to three new editions, and the new one has sold seven thousand before publication.>

<Financially, all to the good, eh?>

<Oh, yes—but frankly I don't know whether these artificial sales do very much good to an author's reputation in the long run. Up like a rocket, down like the stick, you know. When Miss~Vane is released\longdash>

<I am glad you say 'when.'>

<I am not allowing myself to contemplate any other possibility. But \textit{when} that happens, public interest will be liable to die down very quickly. I am, of course, securing the most advantageous contracts I possibly can at the moment, to cover the next three or four books, but I can only really control the advances. The actual receipts will depend on the sales, and that is where I foresee a slump. I am, however, doing well with serial rights, which are important from the point of view of immediate returns.>

<On the whole, as a business man, you are not altogether glad that this has happened?>

<Taking the long view, I am not. Personally, I need not say that I am extremely grieved, and feel quite positive that there is some mistake.>

<That's my idea,> said Wimsey.

<From what I know of your lordship, I may say that your interest and assistance are the best stroke of luck Miss~Vane could have had.>

<Oh, thanks—thanks very much, I say—this arsenic book—you couldn't let me have a squint at it, I suppose?>

<Certainly, if it would help you.> He touched a bell. <Miss~Warburton, bring me a set of galleys of <\textit{Death in the Pot}.> Trufoot's are pushing publication on as fast as possible. The book was still unfinished when the arrest took place. With rare energy and courage, Miss~Vane has put the finishing touches and corrected the proofs herself. Of course, everything had to go through the hands of the prison authorities. However, we were anxious to conceal nothing. She certainly knows all about arsenic, poor girl. These are complete, are they, Miss~Warburton? Here you are. Is there anything else?>

<Only one thing. What do you think of Messrs. Grimsby \& Cole?>

<I never contemplate them,> said Mr~Challoner. <Not thinking of doing anything with them, are you, Lord~Peter?>

<Well, I don't know that I am—seriously.>

<If you do, read your contract carefully. I won't say, bring it to us\longdash>

<If ever I do publish with Grimsby \& Cole,> said Lord~Peter, <I'll promise to do it through you.>